\section{Bilateral remittance estimation} \label{sec:remittance-estimation}

To estimate bilaterance remittance flows from U.S. states to Mexican municipalities, we construct a migration-based weighting matrix using annual MCAS from 2010 to 2024. Because these records capture new consular issuances rather than the complete stock of the remitting population, which also includes previously arrived migrants, we rely on the assumption of persistance of migration patterns over time. To mitigate the volatility of yearly migration fluctuation and construct a robust proxy for the long-term spatial distribution of Mexican migrants, we compute the average of these matrices across all years. We then apply this averaged weighting matrix to the state-level remittance outflows from Banxico's series CE168 to derive the U.S. to Mexican municipality bilateral flows. 

\par

Formally, we define $MCAS_{ms}$ as the average annual count of MCAS issuances to migrants from Mexican municipality $m$ residing in U.S. state $s$. For each U.S. state, we calculate the average migration share, $w_{ms}$, which represents the proportion of that state $s$'s Mexican population originating from the $m$ municipality, as the ratio of $MCAS_{ms}$ to the total number of MCAS issuances from that state across all municipalities $M$:
\begin{equation} \label{eq:weights}
    w_{ms} = \frac{MCAS_{ms}}{\sum_{m=1}^{M} MCAS_{ms}}
\end{equation}
By applying the weights computed from equation \eqref{eq:weights} to the official aggregate state-level outflows from Banxico ($r_s$), we derive the estimated bilateral remittance flow ($\hat{r}_{ms}$) from state $s$ to municipality $m$:
\begin{equation}
    \hat{r}_{ms} = r_s \cdot w_{ms}
\end{equation}
This migration share estimation approach yields the final estimate of our outcome variable. 

\section{Underlying assumptions and validation exercises}

The validity of this estimation relies on two primary identification assumptions. First, we assume the temporal persistence of migration flows, implying that the distribution of migrants recorded in the 2010--2024 MCAS data is representative of the active remitting population. Second, we assume that, within a given U.S. state, the average amount remitted does not vary systematically by the migrant's municipality of origin. While this assumption ignores micro-level heterogeneity, it provides a baseline estimate of the bilateral remittance flows that can be validated against the aggregate data provided by Banxico' series CE166. 

\subsection{Persistence of migration networks}

The persistence assumption of migration networks is critical to our estimation strategy, as it justifies the use of MCAS data to construct the migration-based weighting matrix in order to allocate remittance flows. The stability of migration patterns over time is well-documented in the literature, with studies showing that enstablished migration networks tend to dictate future migration patterns, creating persistent flows between specific origin and destination pairs \parencite{card2001, munshi2003}.

\par

To verify this assumption within our specific sample period, we conduct a validation exercise using the MCAS data by computing the Pearsoncorrelation of our constructed weighting matrices across years. As shown in Figure \ref{fig:persistence}, when taking any year from 2012 onward as the base year, the correlation with subsequent years is consistently above 0.90. While the 2010 and 2011 matrices exhibit strong correlation between them, their correlation with later years is smaller, specifically between 0.75 and 0.85. This may be due to structural adjustment in migration patterns following the Great Recession.
\begin{figure}[H]
    \centering
    \IfFileExists{persistence.jpeg}{%
        \includegraphics[width=0.85\textwidth]{persistence.jpeg}%
    }{%
        \fbox{\parbox{0.85\textwidth}{\centering\vspace{2cm}[Figure pending]\vspace{2cm}}}%
    }
    \caption{Correlation of MCAS-based weighting matrices across years. Each pointrepresents the correlation coefficient between the weighting matrix of the base year and the weighting matrix of the comparison year.}
    \caption*{Source: Authors' own elaboration based on MCAS data.}
    \vspace{-16pt}
    \label{fig:persistence}
\end{figure}

This high degree of correlation suggests that the spatial distribution of Mexican migrants across U.S. states has remained highly stable over time, lending credibility to our use of an averaged weighting matrix for estimating bilateral remittance flows. 

\subsection{Validation of the estimated remittance flows}

Our methodology to estimate bilateral remittance flows implicitly assumes a homogeneous send behavior among migants withing a given U.S. destination. That is, migrants in U.S. state $s$ send on average the same amount of remittances, regardless of their specific Mexican municipality of origin $m$. To evaluate the empirical validity of this assumption, we aggregate our estimated bilateral municipal inflows and compare them against the observed municipal receipts reported in Banco de México's series CE166. 

\par 

Figure \ref{fig:estimate-validation} plots the Pearson correlation between our estimates aggregated at the municipal level and the official inflows for each year. The results report correlation coefficients consistently between 0.75 and 0.85,suggesting that our estimation procedure successfully captures a significant portion of the variation in remittance inflows across Mexican municipalities. While this result validates our baseline approach, we acknowledge that it does not capture all the heterogeneity in remittance flows, which is expected given the simplifying assumptions we made.
\begin{figure}[H]
    \centering
    \IfFileExists{estimate_validation.png}{%
        \includegraphics[width=0.85\textwidth]{estimate_validation.png}%
    }{%
        \fbox{\parbox{0.85\textwidth}{\centering\vspace{2cm}[Figure pending]\vspace{2cm}}}%
    }
    \caption{Correlation of estimated municipal remittance inflows with official data. Each point represents the Pearson correlation coefficient between the estimated municipal inflows (aggregated from our bilateral estimates) and the actual municipal inflows recorded in Banxico's series CE166 for a given year.}
    \caption*{Source: Authors' own elaboration.}
    \vspace{-16pt}
    \label{fig:estimate-validation}
\end{figure}
